% O1 Team Contract rubric
\documentclass[11pt]{article}
\usepackage{mae2250-handout}

% --- Rubric metadata ---
\renewcommand{\HandoutType}{Rubric}
\renewcommand{\HandoutNumber}{3}
\renewcommand{\HandoutTitle}{Parameter-Driven Design}
% \renewcommand{\DueDate}{Midnight, your Lab Section day, week of Feb. \nth{9}, and midnight two days after (report)}
% \renewcommand{\TeamType}{Individual}
\usepackage[skip=10pt plus1pt, indent=0pt]{parskip}


\begin{document}
\MakeHandoutHeader

% \section*{Student-facing Rubric}


\subsection*{Student-facing rubric}
\PointsTableAuto{
  \PointsRow{All CAD sketches, bodies, and functional parameters have a meaningful name}{3}
  \PointsRow{All CAD sketches are fully constrained}{1}
  \PointsRow{All parts are assigned a material, and non-driven dimensions are reasonable for that material}{2}
  \PointsRow{The drawn sketch/equations clearly convey design intent and matches the CAD history}{4}
  \PointsRow{Cup would fit in a standard car cup-holder}{1}
  \PointsRow{Cup capacity can be changed with a single parameter without breaking any part}{4}
}
\subsection*{Detailed rubric}
\textbf{Sketches, bodies, and functional parameters have a meaningful name: 3}
\\
\Dock one point for each category of $\{\text{sketches, bodies}\}$ where something is unnamed.
For the parameters, all \emph{driven} and \emph{driving} parameters should have a name (as illustrated in their illustration). If not, \dock 1 point. (Open \texttt{Change parameters} from the \texttt{MODIFY} dropdown menu.)
\par
\textbf{All CAD sketches are fully constrained: 1}
\\
Self explanatory.
\par
\textbf{All parts are assigned a material, with reasonable dimensions: 2}
\\
\Award 1 point for each. When thinking about reasonable dimensions, consider the fabrication process.
\par
\textbf{The drawn sketch/equations clearly convey design intent: 4}
\\
\Award 1 points if the drawn sketch clearly shows which parameters will be driving/driven.\\
\Award 2 point if the sketch and the CAD history match. For example, if the bottom and top of the cup are defined by \texttt{circle sketches} on \texttt{offset planes} and a \texttt{loft} between them, I would expect the drawn sketch to specify the diameters and distance between those planes. If instead the CAD uses an \texttt{extrude} with a \texttt{taper angle}, I would expect that angle to be illustrated in the drawn sketch.
\par
\textbf{Cup would fit in a standard car cup-holder: 1}
\\
\Award 1 point if the bottom 2 inches of the cup are less than 3 inches in diameter.
\par
\textbf{Cup capacity can be changed with a single parameter without breaking any part: 4}
\Award 1 point if the cup total volume can be easily changed between 4 oz and 24 oz, and how to change this is clearly explained through the sketch or other means.
\\
\Award 1 points if the cap always fits correctly within this range.
\\
\Award 1 point if the cup would always fit the car cup holder when changed in these ranges.
\\
\Award 1 point if the aesthetics of the cup remain consistent / appealing when changing the volume in this range.
\\
\end{document}
